\subsection{The uniform signed quotient and the full congruence correction}

\begin{statement}{E.53}{Proved}{The integral common source with the complete double cover}
Retain E.52's actual maps $\rho_m,\rho_R,\Phi_m,\Phi_R$,
where $R=\prod_{p\geq5}\mathbf Z_p$, and retain
$\rho_L:\widehat B\twoheadrightarrow G_L$ from E.38--E.43.
Write $Q_T=\Phi_T(1,V_3(T))\subset C_T$ for
$T=\mathbf Z/m$ or $R$.
For $m>1$ prime to six, the actual modular map
\[
 q_{4m}:\widehat B\twoheadrightarrow
       Q_{4m}=\operatorname{SL}_2(\mathbf Z/4m)/\{\pm I\},
 \quad
 s\longmapsto\left[\begin{pmatrix}0&1\\-1&0\end{pmatrix}\right],
 \quad
 r\longmapsto\left[\begin{pmatrix}0&-1\\1&-1\end{pmatrix}\right]
 \tag{E.53.1}
\]
factors through $\rho_m$. Together with the ordinary
exponent quotient of E.52 it gives the actual isomorphism
\[
 G_m/Q_m\xrightarrow{\sim}\mathcal J_m,
 \qquad
 \mathcal J_m=A_{2m}\mathbin{\times_{A_2}}Q_{4m},
 \qquad |\mathcal J_m|=24m^2|\operatorname{SL}_2(\mathbf Z/m)|.
 \tag{E.53.2}
\]
The two maps to $A_2$ are reduction of exponent sums
modulo two and projective matrix reduction modulo four.
Their values on $s,r$ are the same displayed source values.
Every comparison in \textup{(E.53.2)} commutes with
reduction of $m$.

Put $\mathcal J_1=A_2$, and retain the E.50 epimorphism
$G_L\twoheadrightarrow\mathcal J_5$.
The maximum common quotient for the specified maps
$(\rho_L,\rho_R)$ is exactly $\mathcal J_5$.
For finite $m$ it is $\mathcal J_{\gcd(m,5)}$.
Equivalently their actual joint images are
\[
 \begin{split}
 \operatorname{im}(\rho_L,\rho_R)
       &=G_L\mathbin{\times_{\mathcal J_5}}G_R,\\
 \operatorname{im}(\rho_L,\rho_m)
       &=G_L\mathbin{\times_{\mathcal J_{\gcd(m,5)}}}G_m.
 \end{split}
 \tag{E.53.3}
\]
These are equalities of the source images under the
displayed projections.

Let $\Gamma_5(1)=\ker(\operatorname{SL}_2(\mathbf Z_5)
\to\operatorname{SL}_2(\mathbf F_5))$.
The exact correcting-subgroup image is
\[
 \begin{split}
 \rho_R(R_{\mathrm{nil}}\cap\ker\rho_L)
   =\Phi_R\left(
       \Gamma_5(1)\times\prod_{p\geq7}\operatorname{SL}_2(\mathbf Z_p),
                           V_3(R)\right)
   =:\mathcal K_R.
 \end{split}
 \tag{E.53.4}
\]
Intersecting the source on the left with
$\ker\rho_{D_{23}}$ leaves that image unchanged.
Intersecting additionally with $\ker\rho_5\cap\ker\rho_7$
gives precisely
\[
 \begin{split}
 \mathcal K_R^{35}
   =\Phi_R\left(
       \Gamma_5(1)\times\Gamma_7(1)
          \times\prod_{p\geq11}\operatorname{SL}_2(\mathbf Z_p),
          \ 5V_3(\mathbf Z_5)\times7V_3(\mathbf Z_7)
               \times\prod_{p\geq11}V_3(\mathbf Z_p)\right).
 \end{split}
 \tag{E.53.5}
\]
Here the argument of $\Phi_R$ is the displayed element
of $\operatorname{SL}_2(R)\times V_3(R)$.
Each prescribed target value in either correcting image
has the full source fibre $k\ker(\rho_R|_K)$, where
$K$ is the corresponding displayed intersection and
$k\in K$ is one preimage.
Every such correction maps to the identity in every
nilpotent continuous quotient of $\widehat H$.
\end{statement}

\begin{proof}
We retain the central sign first, and then determine
the kernel of each actual joint source map.

\paragraph{The canonical signed matrix on the pure subgroup.}
For a pure word $f\in F$, its ordinary integral modular
matrix is determined up to its central sign. Exactly
one sign has both diagonal entries equal to one modulo
four. Indeed the off-diagonal entries are even, and
determinant one makes the two odd diagonal residues
equal modulo four. Products of matrices with diagonal
entries one modulo four retain that property.
The sign choice therefore defines a homomorphism on
the pure subgroup, with
\[
 x\longmapsto\begin{pmatrix}1&2\\0&1\end{pmatrix},\qquad
 y\longmapsto\begin{pmatrix}1&0\\-2&1\end{pmatrix}.
 \tag{E.53.6}
\]
For finite quotients this means the compatible CRT
representative with the indicated component modulo four;
it also defines the corresponding continuous pure map.
With $P=\left(\begin{smallmatrix}0&1\\1&0\end{smallmatrix}\right)$,
its component modulo $m$ is $P\theta(f)P$.
For exponent sums $(a,b)$ its component modulo four is
\[
                   \begin{pmatrix}1&2a\\-2b&1\end{pmatrix}.
 \tag{E.53.7}
\]
Both assertions follow on the two free generators
from \textup{(E.53.6)} and extend by multiplication.
In particular $Q_4\simeq A_2$ through this exact map.

An element of $\ker\rho_m$ lies in $K_2$ by E.52's
two parity characters and has $\theta$-value one.
Equations \textup{(E.53.6)--(E.53.7)} make its
canonical matrix the identity modulo $4m$.
Thus $q_{4m}$ factors through $G_m$.
It is onto: the product of the displayed $s,r$ matrices
is $E_{12}(-1)$; its powers and conjugation by $s$
give all upper and lower elementary matrices.
Elementary elimination over each prime-power factor
and CRT prove generation of the full determinant-one
group, hence of its specified projective quotient.

The kernel of $G_m\to A_{2m}$ is $C_m$.
On $C_m$, the modular map has trivial four-component
and its signed $m$-component is exactly $P\theta P$.
The common kernel is consequently $Q_m$.
The image of $G_m/Q_m$ lies in the actual fibre
product in \textup{(E.53.2)}.
Now $|Q_{4m}|=24|\operatorname{SL}_2(\mathbf Z/m)|$,
because $|\operatorname{SL}_2(\mathbf Z/4)|=48$ and
the diagonal central sign has order two.
The fibre product therefore has the order stated in
\textup{(E.53.2)}, equal to $|G_m/Q_m|$ by E.52.
The injective map is onto. This proves both the
signed comparison and every transition map.

\paragraph{Finite quotients of the integral $\operatorname{SL}_2$ factor.}
We give the group facts needed to control the full
inverse system, including its central quotients.
For $p\geq5$ put
$L_p=\operatorname{SL}_2(\mathbf Z_p)$ and
$\Gamma_p(n)=\ker(L_p\to\operatorname{SL}_2(\mathbf Z/p^n))$.
The map
\[
 \Gamma_p(n)/\Gamma_p(n+1)\longrightarrow\mathfrak{sl}_2(\mathbf F_p),
              \qquad 1+p^n A\longmapsto A\bmod p
 \tag{E.53.8}
\]
is an isomorphism of additive groups and is equivariant
for the actual conjugation action of $L_p$.
The determinant equation gives trace zero; the two
elementary matrices and $\operatorname{diag}(1+p^na,(1+p^na)^{-1})$
give every basis vector, proving surjectivity.

This adjoint module is irreducible.
Conjugation by $E_{12}(t)$ and $E_{21}(t)$ is a
polynomial of degree at most two in $t$.
Interpolation at $0,1,-1$ shows that an invariant
subspace is stable under brackets with the usual
$e,f$, and hence with $h=[e,f]$.
The $h$-weights $0,2,-2$ are distinct for $p\geq5$.
Projection onto those three eigenspaces, followed
by brackets with $e,f$, shows that any nonzero
invariant subspace contains $e,f,h$.

Let $U\triangleleft L_p$ be open.
The group $\operatorname{PSL}_2(\mathbf F_p)$ is simple
by the following concrete argument. Its faithful
projective action is primitive; the translations form
an abelian normal subgroup of a point stabilizer,
and their conjugates generate the group. The elementary
commutators E.52.17 make it perfect.
A nontrivial normal subgroup is transitive in this
primitive faithful action. Modulo that subgroup,
all conjugate translation groups have the same
abelian image, because every group element is a
product of a normal-subgroup element and a point
stabilizer. The quotient is abelian, and perfectness
makes it trivial. This proves simplicity.
Consequently the reduction of $U$ modulo $p$ is
central or is the whole group. In the latter case
$L_p/U$ is a finite $p$-group. It is also perfect
by E.52.17, so it is trivial.
The full-reduction conclusion also holds for a closed normal
subgroup: its quotient is pro-$p$ and topologically perfect;
every finite quotient is a perfect finite $p$-group and is
therefore trivial.

In the former case consider $U\cap\Gamma_p(1)$.
Let $n$ be the first layer in \textup{(E.53.8)}
where it has nonzero image. Irreducibility makes
that image the whole layer. Choose three elements
of $U$ giving a basis there.
For $a\geq1$, direct binomial expansion gives
\[
           (1+p^a A)^p\equiv1+p^{a+1}A\pmod{p^{a+2}}.
 \tag{E.53.9}
\]
Successive powers of the chosen elements therefore
span every later layer. Successive approximation
and closedness prove
$U\cap\Gamma_p(1)=\Gamma_p(n)$.
The reduction of $U$ is either $1$ or $\{\pm I\}$.
In the second case write an element above $-I$ as
$-v$ with $v\in\Gamma_p(1)$.
Its square lies in $\Gamma_p(n)$.
The finite $p$-group $\Gamma_p(1)/\Gamma_p(n)$
has no element of order two, so $v\in\Gamma_p(n)$.
Thus the possibilities are precisely
\[
                  U=\Gamma_p(n),\quad
                  U=\{\pm I\}\Gamma_p(n),\quad U=L_p.
 \tag{E.53.10}
\]
The argument also proves that any closed normal
subgroup with a nonzero first congruence layer
contains all of $\Gamma_p(1)$.

Every nontrivial finite quotient of $L_p$ therefore
has $\operatorname{PSL}_2(\mathbf F_p)$ as a
composition factor, and its centre has exponent
dividing two. For the last assertion a matrix
commuting with the two elementary generators is
scalar, and the determinant equation gives $\pm I$.
For a quotient by the central sign, reduce first
modulo $p$. A central projective element is scalar
there; its commutators with the elementary generators
are consequently both signs with reduction $+1$,
so they equal $+1$ before taking the quotient.
The same scalar calculation finishes that case.

A continuous finite quotient of $\prod_{p\geq5}L_p$
uses finitely many factors. Their images commute
and generate it. If a product of elements from
these factor images is central, each factor element
commutes with its own factor image, because all
other factor images commute with that image already.
Each factor element has square one by the preceding
paragraph, so the whole central product has square
one. We have proved that every such finite quotient
has centre of exponent dividing two.

\paragraph{The exact lower-central correcting image.}
By E.43, $\rho_L(R_{\mathrm{nil}})=C_L$ is perfect.
By E.52, $\rho_R(R_{\mathrm{nil}})=C_R$.
Apply the profinite Goursat argument to their actual
joint image and put
$\mathcal K_R=\rho_R(R_{\mathrm{nil}}\cap\ker\rho_L)$.
Then $C_R/\mathcal K_R$ is a finite quotient of
$C_L$, hence perfect.
The additive factor $Q_R$ is central in $C_R$.
In this perfect quotient the image of
$\operatorname{SL}_2(R)$ is the whole group: the
commutator of the product of that image and the
central additive image is just the first image.
The additive image is consequently in the centre
of a finite quotient of $\operatorname{SL}_2(R)$.
That centre has exponent dividing two, while every
finite quotient of $Q_R$ has order prime to six.
Thus $Q_R\subset\mathcal K_R$.

Let $U$ be the resulting closed normal subgroup
of $\operatorname{SL}_2(R)$.
The literal E.38 source order is
$|G_L|=2^{82}3^{26}5^{21}7^2$.
For $p\geq11$, a nontrivial quotient of $L_p$
would have the factor $\operatorname{PSL}_2(\mathbf F_p)$,
whose order is divisible by $p$.
Its order cannot divide $|G_L|$.
Thus the projection of $U$ onto every such $L_p$
is all of $L_p$.
For $p=7$, E.51 proves that the very same
$R_{\mathrm{nil}}\cap\ker\rho_L$ maps onto
$\operatorname{SL}_2(\mathbf F_7)$.
Its closed normal projection in $L_7$ is therefore
all of $L_7$, by the full-reduction case above.
For any one of these primes, commutating elements
of $U$ with elements supported in that one factor
shows that $U$ contains the pure factor $L_p$,
because $L_p$ is topologically perfect.
Closedness then gives the entire product for $p\geq7$.

The five-component is determined by an actual word.
In the complete signed Ferrari source, both
$\rho_L(x)$ and $\rho_L(y)$ have exact order $30$.
In particular $x^{60},y^{60}$ belong to
$\ker\rho_L$.
Choose $c\in R_{\mathrm{nil}}$ whose five-component
under $\theta\rho_R$ is $E_{12}(1)$, using E.52's
surjectivity of this same closed subgroup.
Then $[x^{60},c]$ lies in
$R_{\mathrm{nil}}\cap\ker\rho_L$, since the residual
is normal in $\widehat B$.
Its five-component is the commutator of $E_{21}(120)$
and $E_{12}(1)$, and modulo $25$ it is
\[
                  \begin{pmatrix}6&20\\0&21\end{pmatrix}.
 \tag{E.53.11}
\]
Its nonzero first-layer vector is
$\left(\begin{smallmatrix}1&4\\0&4\end{smallmatrix}\right)$.
The three explicit conjugate vectors in the
certificate span $\mathfrak{sl}_2(\mathbf F_5)$.
Equation \textup{(E.53.9)} and closedness give the
whole $\Gamma_5(1)$ in the five-projection of $U$.
The other factors already belong to $U$ separately,
so $\Gamma_5(1)$ belongs as a pure factor as well.
Conversely E.50's actual common signed quotient
gives the map $C_L\twoheadrightarrow\operatorname{SL}_2(\mathbf F_5)$
whose composition with the source is $\theta\rho_5$.
It forces the five-reduction of every element of
$U$ to be the identity.
This proves precisely \textup{(E.53.4)}.

\paragraph{The whole common source.}
Put $M=\rho_R(\ker\rho_L)$.
The two maps to $\mathcal J_5$ agree by E.50 and
\textup{(E.53.2)}, so
$M\subset\ker(G_R\to\mathcal J_5)$.
It contains $\mathcal K_R$.
Modulo $\mathcal K_R$, the restriction of the latter
map to $C_R$ is the isomorphism onto its
$\operatorname{SL}_2(\mathbf F_5)$ subgroup.
Consequently its remaining kernel injects into
the pure exponent quotient, and is exactly the
kernel of reduction of the two even exponent
coordinates from $R$ to $\mathbf F_5$.
Surjectivity onto that exponent kernel follows by
correcting any chosen lift by its uniquely required
$\operatorname{SL}_2(\mathbf F_5)$ component.

The images of $x^{60},y^{60}$ lie in $M$, have
trivial signed component modulo five and have
even exponent coordinates $(30,0),(0,30)$.
Their closed spans in $R^2$ are exactly $5R^2$,
because six is a unit in $R$ and ordinary integers
are dense in $R$ through CRT.
They therefore generate the entire remaining
kernel just described. It follows that
\[
                       M=\ker(G_R\to\mathcal J_5).
 \tag{E.53.12}
\]
The actual joint source is a subdirect product;
\textup{(E.53.12)} and its agreed quotient identify
it with the full fibre product in
\textup{(E.53.3)}.
Reduction to any finite $m$ gives the second
identity: only a first factor of five can remain
in the displayed correcting kernel and exponent
kernel; with no factor of five the quotient is
exactly $A_2$.

\paragraph{Retaining the old small-prime and symplectic values.}
Every finite solvable quotient of $\mathcal K_R$
has order prime to six. Its perfect $L_p$ factors
have trivial solvable image, and the remaining
factors are the pro-five $\Gamma_5(1)$ and the
additive product at primes at least five.
The group $D_{23}$ is a finite solvable group whose
order has only primes two and three.
Their common quotient is therefore trivial.
Applying Goursat to the actual image of
$R_{\mathrm{nil}}\cap\ker\rho_L$ proves that
intersecting its source with $\ker\rho_{D_{23}}$
still maps onto all of $\mathcal K_R$.

The maps $\rho_5,\rho_7$ factor through $\rho_R$.
Hence imposing their kernels on a source mapping
onto $\mathcal K_R$ gives exactly the intersection
of $\mathcal K_R$ with their two target kernels.
In the coordinates of E.52 that is precisely
\textup{(E.53.5)}.
The same argument about finite solvable quotients
applies after this intersection, since the added
seven-component is pro-seven.
Finally, a homomorphism from $\widehat H$ to a
nilpotent profinite group of finite nilpotency
class kills some $\gamma_j\widehat H$ and hence
$R_{\mathrm{nil}}$. For an inverse limit of such
quotients it kills this subgroup at every stage,
and so kills it in the limit.
The full coset assertions follow from the kernels
of these specified continuous homomorphisms.
\end{proof}

The exact program
\texttt{e\_symplectic\_integral\_gl\_comparison\_root.py}
reconstructs the complete signed source without executing
the older distribution script that writes an output at
import time. It checks the original signed model arrays,
the exact orders of $x,y$ in all 36 induced coordinates,
all 102 canonical signed modular comparisons, the
level-25 commutator and its three independent conjugate
Lie vectors. Its run takes 3.779 seconds.
The first run used an overstrong exact-order assertion
of 60, which the computation rejected. That script and
traceback are preserved. The actual order is 30; both
power relations used in the proof are verified in the
complete source, and no conclusion uses the rejected
exact-order equality.

An independent complete proof audit and isolated replay reproduce
every mathematical JSON entry in 4.551 seconds. A separate
symbolic commutator calculation gives the same level-25
matrix, and the three conjugate vectors have determinant
$3$. All protected originals are unchanged.
